% =====================================================================
%  THÈME 4 — Logarithmes, pH & incertitude — EXERCICES — Niveau MOYEN
% =====================================================================
\documentclass[11pt,a4paper]{article}
\usepackage[utf8]{inputenc}
\usepackage[T1]{fontenc}
\usepackage{lmodern}
\usepackage[french]{babel}
\usepackage{amsmath,amssymb,mathtools}
\usepackage{icomma}
\usepackage{siunitx}
\sisetup{output-decimal-marker={,},per-mode=symbol}
\usepackage[version=4]{mhchem}
\usepackage{xcolor}
\usepackage[most]{tcolorbox}
\usepackage{enumitem}
\usepackage{booktabs}
\usepackage{fancyhdr}
\usepackage[a4paper,margin=2.2cm,headheight=26pt,headsep=8pt]{geometry}

\definecolor{primary}{RGB}{150,98,162}
\definecolor{primarydark}{RGB}{92,52,110}
\newcommand{\cs}{chiffre significatif}
\newcommand{\CS}{chiffres significatifs}

\pagestyle{fancy}\fancyhf{}
\fancyhead[L]{\small\textcolor{primarydark}{\textbf{Fiche 1} — Chiffres significatifs}}
\fancyhead[R]{\small\textcolor{primarydark}{\textsc{Exercices · Moyen · Thème 4}}}
\fancyfoot[C]{\small\thepage}
\renewcommand{\headrulewidth}{0.8pt}
\renewcommand{\headrule}{\hbox to\headwidth{\color{primary}\leaders\hrule height \headrulewidth\hfill}}

\newtcolorbox[auto counter]{exo}[1][]{enhanced,breakable,boxrule=0.8pt,arc=2pt,
  colback=primary!4,colframe=primary,left=3mm,right=3mm,top=2mm,bottom=2mm,
  fonttitle=\bfseries,coltitle=white,
  title={Exercice~\thetcbcounter\ifx\relax#1\relax\relax\else\ ---\ \textmd{\itshape #1}\fi}}

\setlength{\parindent}{0pt}\setlength{\parskip}{3pt}

\begin{document}

\begin{center}
{\Large\bfseries\color{primarydark} Exercices — Niveau Moyen}\\[1pt]
{\color{primary} Thème 4 : calculs de pH \& incertitude relative}
\end{center}
\vspace{1mm}
\textit{Les valeurs de logarithme utiles sont fournies dans les énoncés (aucune calculatrice
nécessaire).}
\vspace{2mm}

\begin{exo}[Calculer un pH avec le bon nombre de décimales]
On donne : $\log(\num{2,5})=\num{0,398}$ ; $\log(\num{3,0})=\num{0,477}$ ; $\log(\num{8,00})=\num{0,903}$.
Calcule le pH avec le bon nombre de décimales.
\begin{enumerate}[label=\alph*),leftmargin=1.6em,itemsep=1pt]
  \item $[\ce{H3O+}]=\num{2,5e-4}\,\si{\mol\per\litre}$
  \item $[\ce{H3O+}]=\num{3,0e-9}\,\si{\mol\per\litre}$
  \item $[\ce{H3O+}]=\num{8,00e-6}\,\si{\mol\per\litre}$
\end{enumerate}
\end{exo}

\begin{exo}[Remonter du pH à la concentration]
On donne : $10^{\num{0,40}}=\num{2,5}$ et $10^{\num{0,48}}=\num{3,0}$. Détermine $[\ce{H3O+}]$ avec le
bon nombre de \CS.
\begin{enumerate}[label=\alph*),leftmargin=1.6em,itemsep=1pt]
  \item $\mathrm{pH}=\num{3,60}$ \qquad b) $\mathrm{pH}=\num{8,52}$
\end{enumerate}
\end{exo}

\begin{exo}[Incertitude relative d'une pesée]
Calcule l'incertitude relative $\dfrac{\Delta x}{x}$ (en \%) pour chaque mesure.
\begin{enumerate}[label=\alph*),leftmargin=1.6em,itemsep=1pt]
  \item $m = \qty{13,02\pm0,01}{\gram}$ \qquad b) $V = \qty{20,96\pm0,01}{\milli\litre}$
\end{enumerate}
\end{exo}

\begin{exo}[Lire la précision dans le nombre de \CS]
Trois mesures d'une même longueur sont rapportées : $\qty{1,2e2}{\milli\metre}$ ($2$~CS),
$\qty{1,20e2}{\milli\metre}$ ($3$~CS) et $\qty{1,200e2}{\milli\metre}$ ($4$~CS).
\begin{enumerate}[label=\alph*),leftmargin=1.6em,itemsep=1pt]
  \item Donne l'ordre de grandeur de l'incertitude \emph{relative} de chacune (règle $10^{-(n-1)}$).
  \item Laquelle est la plus précise, et pourquoi ?
\end{enumerate}
\end{exo}

\begin{exo}[Analyser une série de mesures]
La vraie valeur d'un volume à doser est $\qty{25,00}{\milli\litre}$. Trois binômes répètent la mesure
$4$ fois :
\begin{center}\small\renewcommand{\arraystretch}{1.15}
\begin{tabular}{c l}
\toprule
\textbf{Binôme} & \textbf{Mesures (\si{\milli\litre})}\\
\midrule
A & $\num{24,98}$ ; $\num{25,01}$ ; $\num{24,99}$ ; $\num{25,02}$\\
B & $\num{25,50}$ ; $\num{25,52}$ ; $\num{25,49}$ ; $\num{25,51}$\\
C & $\num{24,60}$ ; $\num{25,40}$ ; $\num{25,10}$ ; $\num{24,90}$\\
\bottomrule
\end{tabular}
\end{center}
Pour chaque binôme, dis si les mesures sont \emph{justes}, \emph{fidèles}, les deux, ou aucune.
Justifie brièvement.
\end{exo}

\end{document}
